\begin{example}[Для не T1-пространств из limit point compactness не следует счётная компактность.]
    Пусть $X = \N$ и базой топологии является $\beta = \left\{[2n - 1, 2n] \mid n \in \N\right\}$.
    Очевидно, что пространство не удовлетворяет первой аксиоме отделимости -- точки не замкнуты.

    Такое пространство не является счётным компактом -- в качестве покрытия можно взять само $\beta$ и из него, очевидно, нельзя выделить конечное подпокрытие.

    Но оно является limit point compact.
    Если $E \subset \N$ -- бесконечное множество, то для некоторого $k$ или $2k - 1$, или $2k$ лежит в множестве $E$.
    Но если $2k$ (или $2k - 1$) лежит в $E$, то $2k - 1$ ($2k$, соответственно) является предельной точкой $E$: любая окрестность $U(2k - 1)$ (или $U(2k)$) содержит вторую точку из пары, а, следовательно, пересечение с $E$ будет непусто.
\end{example}
\begin{proposition}
    Если $(X, \tau) \in \cat{Top}$ и $X$ удовлетворяет первой аксиоме отделимости, то если непустое $K \subset X$ -- limit point compact, то $K$ является и счётным компактом.
\end{proposition}
\begin{proof}
    Пусть $K$ не является счётным компактом, то есть существует такое счётное покрытие $P = \{U_n\}_{n = 1}^{\infty} \subset \tau$, что из него нельзя выделить конечное подпокрытие.
    А значит $\forall m \in X \exists x_m \in K \setminus \bigcup_{k = 1}^{m} U_k$.
    Рассмотрим $E\colon$
    \[
        E = \left\{x_m \mid m \in \N\right\}.
    \]
    Оно имеет предельную точку (поскольку бесконечно) $x$, которая лежит в некотором $U_{m_0} \in P$ -- поскольку $P$ является покрытием.
    В $U_{m_0}$ не лежит никакой $x_m$ при $m \geq m_0$ -- по построению.
    Пусть теперь $V = U_{m_0} \setminus \bigcup_{k = 1}^{m_0 - 1} \{x_k\}$.
    Оно является открытым в силу замкнутости всех точек в $X$.
    В силу предельности $x$ для $E$ существует $x_m \in V \setminus x \cap E$.
    Но $m$ не может быть меньше $x_{m_0}$ -- мы ручками их все выкинули.
    А следовательно $m \geq m_0$ -- противоречие, они тоже не лежат в $V$ по построению.

    Следовательно, $K$ является счётным компактом.
\end{proof}
\section{Центрированные системы множеств и критерий компактности.}
\begin{definition}
    Пусть $(X, \tau) \in \cat{Top}$ и $S \subset X$ -- некоторое подмножество $X$.
    Определим топологию подпространства (индуцированную) на $S$ как
    \[
        \tau_S = \left\{U \cap S \mid U \in \tau\right\}.
    \]
\end{definition}
\begin{definition}
    Если $(X, \tau) \in \cat{Top}$ и $S \subset X$, то будем говорить что $W \in \tau_S$ -- открыто в $S$.

    Аналогично, $S \setminus W$ называется замкнутым в $S$.
\end{definition}
\begin{proposition}
    Пусть $(X, \tau) \in \cat{Top}$.
    Тогда $K \subset X$ -- компактно тогда и только тогда, когда $(K, \tau_K)$ -- компактно.
\end{proposition}
\begin{proof}
    Это очевидно, в силу того, что всякое открытое покрытие $K$ в $X$ будет открытым и в $(K, \tau_K)$.
    Верно и обратное -- если есть некоторое открытое покрытие в $(K, \tau_K)$, то так как для всякого открытого в $K$ множества $S$ существует открытое в $X$ множество $U$ такое, что $S = U \cap K$, то у нас есть и открытое покрытие $K$ в $X$.
\end{proof}
\begin{definition}
    Пусть $X$ -- некоторое множество.
    Семейство множеств $F \subset 2^X$ называется центрированным, если пересечение всякой конечной подсистемы $\{S_i\}_{i = 1}^{N} \subset F$ непусто: $\bigcap_{i = 1}^{N} S_i \neq \emptyset$.
\end{definition}
\begin{proposition}
    Если $(X, \tau) \in \cat{Top}$ и любая центрированная система $F$ замкнутых в $K \subset X$ подмножеств имеет непустое пересечение, то $K$ -- компактно.

    Верно и обратное.
\end{proposition}
\begin{proof}
    Пусть $K$ -- компакт и $F$ -- центрированная система замкнутых в $K$ множеств.
    Пусть $\bigcap_{S \in F} S = \emptyset$.
    Тогда
    \[
        K = K \setminus \bigcap_{S \in F} S = \bigcup_{S \in F} K \setminus S.
    \]
    Поскольку всякое $S$ из $F$ замкнуто в $K$, то $K \setminus S$ -- открыто в $K$ и, следовательно, семейтсво
    \[
        P = \left\{K \setminus S \mid S \in F\right\},
    \]
    является открытым покрытием $K$.
    Но, тогда, в силу компактности существует конечное подпокрытие $\{K \setminus S_i\}_{i = 1}^{N}$.
    То есть
    \[
        K = \bigcup_{i = 1}^{N} K \setminus S_i = K \setminus \bigcap_{i = 1}^{N} S_i.
    \]
    Но пересечение любой конечной подсистемы $F$ -- непусто, а следовательно мы приходим к противоречию.

    Пусть теперь всякая центрированная система $F$ из замкнутых в $K$ множеств имеет непустое пересечение.
    Пусть $P$ -- произвольное открытое покрытие $K$:
    \[
        K = \bigcup_{W \in P} W.
    \]
    Пусть никакое конечное подсемейство $\{W_i\}_{i = 1}^{N}$ не является подпокрытием $K$, то есть
    \[
        \emptyset \neq K \setminus \bigcup_{i = 1}^{N} W_i = \bigcap_{i = 1}^{N} K \setminus W_i.
    \]
    Рассмотрим семейство
    \[
        F = \left\{ K \setminus W \mid W \in P \right\}.
    \]
    Заметим, что оно обладает свойством центрированности, как показано выше, и состоит из замкнутых в $K$ множеств.
    Тогда пересечение всех $\bigcap_{S \in F} S \neq \emptyset$, и, значит
    \[
        \emptyset \neq \bigcap_{W \in P} K \setminus W = K \setminus \bigcup_{W \in P} W.
    \]
    Но тогда $P$ не является покрытием -- противоречие, а значит $K$ -- компактно.
\end{proof}
\begin{corollary}
    Если $(X, \tau)$ -- хаусдорфово пространство и $F$ -- непустое центрированное семейство компактов в $X$, то
    \[
        \bigcap_{K \in F} K \neq \emptyset.
    \]
\end{corollary}
\begin{proof}
    Заметим, что существует $K_0 \in F$ (так как $F$ -- непусто).
    Тогда
    \[
        \left\{S_K = K \cap K_0 \mid K \in F\right\}
    \]
    является центрированной системой множеств, замкнутых в $K_0$ (поскольку компакты $K$ замкнуты в $X$, в силу хаусдорфовости).
    В силу компактности $K_0$ её пересечение непусто, что и завершает доказательство.
\end{proof}
\section{Компактность в метрических пространствах.}
\begin{definition}
    Пусть $(X, \rho)$ -- метрическое пространство.
    Тогда множество $S \subset X$ называется вполне ограниченным, если
    \[
        \forall \epsilon > 0 \ \exists \{y_i\}_{i = 1}^{N} \subset S\colon S \subset \bigcup_{i = 1}^{N} O_{\epsilon}(y_i).
    \]
    Множестве $\{y_i\}_{i = 1}^{N}$ называется $\epsilon$-сетью.
\end{definition}
\begin{theorem}[Критерий Фреше]
    Для метрического пространства $(X, \rho)$ и $K \subset X$ следующие условия эквивалентны:
    \begin{enumerate}
        \item $K$ -- компакт;
        \item $K$ -- полно и вполне ограничено;
        \item $K$ -- секвенциальный компакт.
    \end{enumerate}
\end{theorem}
\begin{proof}
    Пусть $K$ -- компактно.
    Тогда семейство шаров
    \[
        P = \left\{O_\epsilon(x) \mid x \in K\right\},
    \]
    является открытым покрытием $K$ и, следовательно, допускает конечное подпокрытие, которое и будет являться искомой эпсилон сетью.
    Таким образом $K$ -- вполне ограничено.

    Заметим, что любая убывающая система непустых замкнутых множеств $\{S_j\}$ будет являться центрированной в $K$ и, следовательно, будет иметь непустое пересечение (в силу критерия компактности через центрированные системы).
    Тогда, если $\{x_n\} \subset K$ -- фундаментальная последовательность, то последовательность $S_k = [\{x_n\}_{n \geq k}]$ будет являться как раз-таки такой системой и
    \[
        \bigcap_{k = 1}^{\infty} S_k \neq \emptyset.
    \]
    Пусть $x \in \bigcap_{k = 1}^{\infty} S_k$.
    Поскольку последовательность $\{x_n\}$ -- фундаментальна, то
    \[
        \forall \epsilon > 0 \exists N(\epsilon) \in \N\colon \forall m, n \geq N(\epsilon) \hookrightarrow \rho(x_m, x_n) < \epsilon.
    \]
    С другой стороны, так как $x$ лежит в замыкании $S_{N(\epsilon)}$, то существует $k \geq N(\epsilon)$ такое, что $\rho(x, x_k) \leq \epsilon$.
    Тогда
    \[
        \forall m \geq N(\epsilon) \hookrightarrow \rho(x, x_m) \leq \rho(x, x_k) + \rho(x_k, x_m) \leq 2\epsilon,
    \]
    и это означает что $x_n \xrightarrow{\rho} x$ -- и, следовательно, $K$ -- полно.

    Пусть теперь $K$ -- полно и вполне ограничено.
    Рассмотрим произвольную последовательность $\{x_n\} \subset K$.
    Пусть $\{y_i\}_{i = 1}^{N}$ -- конечная $1$-сеть для $K$.
    Существует шар $B_1(y_j)$ такой, что в нём лежит бесконечное число элементов $\{x_n\}$.
    Обозначим её за $\{x_{n(1)}\} \subset B_1(y_j)$.
    Аналогично выберем $\frac{1}{2}$-сеть в нём, и выберем в нём шар, содержащий бесконечную подпоследовательность.
    Действуя итеративно, мы получим последовательность $\{x_{n(k)}\}$ такую, что $\forall k, l \geq K \hookrightarrow \rho(x_{n(l)}, x_{n(k)}) \leq \frac{2}{K}$.
    Следовательно, эта подпоследовательность -- фундаментальна и, в силу полноты $K$, сходится к некоторой точке из $K$, что и показывает его секвенциальную компактность.

    \begin{note}
        Заметим, что верно и обратное включение -- что если из всякой последовательности элементов множества можно выделить фундаментальную подпоследовательность, то и множество будет являться вполне ограниченным.

        И это действительно очевидно: предположим противное, то есть что множество не является вполне ограниченным.
        Тогда существует $\epsilon_0$ такое, что не существует конечной $\epsilon_0$-сети.
        Следовательно, существует $\epsilon_0$-разреженная последовательность из элементов $S$ -- но из этой последовательности нельзя выделить фундаментальную подпоследовательность, противоречие.
    \end{note}

    Пусть теперь $K$ -- секвенциальный компакт.
    Покажем вполне ограниченность $K$.
    Она очевидно, поскольку если у нас нет конечной $\epsilon$-сети, то существует счётное $\epsilon$-разреженное множество и, следовательно, $\epsilon$-разреженная последовательность.
    Она не имеет никакой сходящейся подпоследовательности, что вступает в противоречие с секвенциальной компактностью.

    Продолжим доказательство критерия Фреше.
    Пусть $K$ -- секвенциальный, но не топологический компакт.
    Тогда существует открытое покрытие $P$ множества $K$ такое, что оно не имеет конечного подпокрытия.
    Выберем $\forall m \in \N$ конечную $\frac{1}{m}$-сеть для $K$.
    Поскольку не существует конечного подпокрытия $P$, то некоторый шар $O\frac{1}{m}(y_m)$ не может быть покрыт конечным числом элементов покрытия $P$.
    \begin{note}
        Заметим, что если $K$ -- вполне ограничено, то $\forall \epsilon > 0$ существует $\{z_i\} \subset X$ такие, что $K \subset O_\epsilon(z_i)$. % Чё блять?
    \end{note}

    И, тогда, у нас есть последовательность $\{y_m\} \subset K$.
    В силу секвенциальной компактности у нас существует подпоследовательность $\{y_{m_s}\}$ такая, что $y_{m_s} \xrightarrow{\rho} y_0 \in K$.
    Эта точка лежит в некотором $V_0 \in P$ (так как $P$ -- покрытие) и, следовательно, в силу открытости $V_0$, существует $r_0 > 0\colon y_0 \in O_{r_0}(y_0) \subset V_0$.
    Но тогда все элементы последовательности кроме конечного числа попадают в этот шар и, следовательно, у нас есть конечное покрытие элементом $P$ -- $V_0$ -- некоторого шара $O_{1/m_{s_0}}(y_{m_{s_0}})$, что противоречит выбору последовательности.
    Значит $K$ -- компакт.
\end{proof}
\begin{note}
    Если $(X, \rho)$ -- полно и $S \subset X$ -- замкнуто, то $S$ -- полно.

    Это очевидно в силу того, что если у нас есть некоторая фундаментальная последовательность $\{x_n\} \subset S$, то она сходится в $X$ в силу его полноты к некоторому $x^*$.
    Но, тогда, $x^*$ лежит в замыкании $S$ -- которое совпадает с ним самим.
    Значит, всякая фундаментальная последовательность из $S$ сходится в $S$ и, таким образом, оно является полным.
\end{note}
\begin{note}

\end{note}


\section{Сепарабельность.}
\begin{definition}
    Бесконечное метрическое пространство $(X, \rho)$ называется сепарабельным, если существует счётное всюду плотное множество $S \subset X$, то есть такое множество, что
    \[
        [S] = X.
    \]
\end{definition}
\begin{example}
    Если $(K, \rho)$ -- бесконечный метрический компакт, то он является сепарабельным.
\end{example}
\begin{proposition}[Критерий несепарабельности.]
    Для бесконечного метрического пространства $(X, \rho)$ следующие утверждения эквивалентны:
    \begin{enumerate}
        \item $X$ -- несепарабельно;
        \item В $X$ существует более чем счётное $\epsilon$-разреженное множество для некоторого $\epsilon > 0$, то есть $S \subset X\colon$
        \[
            \forall x, y \in S\colon x \neq y \hookrightarrow \rho(x, y) \geq \epsilon.
        \]
    \end{enumerate}
\end{proposition}
\begin{proof}
    Пусть у нас существует более чем счётное $\epsilon$-разреженное множество $S$.
    Пусть $A \subset X$ -- всюду плотно.
    Тогда
    \[
        \forall x \in S \ \exists a(x) \in A\colon \rho(x, a(x)) < \frac{\epsilon}{2}.
    \]
    А значит
    \[
        \forall x, y \in S\colon x \neq y \hookrightarrow \epsilon \leq \rho(x, y) \leq \rho(x, a(x)) + \rho(a(x), a(y)) + \rho(a(y), y) < \epsilon + \rho(a(x), a(y)).
    \]
    Поэтому
    \[
        \rho(a(x), a(y)) > 0.
    \]
    Но тогда мы получаем биекцию между некотором подмножеством $A$ и $S$, а значит $A$ -- более чем счётно и $(X, \rho)$ -- несепарабельно.

    Пусть теперь пространство несепарабельно.
    В нём у нас существуют две различные точки $x_0, y_0 \in X\colon \epsilon_0 = \rho(x_0, y_0) > 0$.
    Теперь $\forall \epsilon \in (0, \epsilon_0)$ будем рассматривать семейства множеств вида
    \[
        \Phi_{\epsilon} = \left\{ A \subset X \mid |A| > 1, \ \forall x, y \in A\colon x \neq y \hookrightarrow \rho(x, y) \geq \epsilon \right\}.
    \]
    Заметим, что $\forall \epsilon \in (0, \epsilon_0)$ семейство $\Phi_\epsilon \neq \emptyset$, так как $\{x_0, y_0\}$ в нём точно лежит.

    Упорядочим $\Phi_\epsilon$ относительно вложения.
    Рассмотрим некоторую цепь $L_\epsilon \subset \Phi_\epsilon$.
    Пусть
    \[
        A_\epsilon = \bigcup_{A \in L_\epsilon} A.
    \]
    Пусть $x, y \in A_\epsilon\colon x \neq y$.
    В силу линейной упорядоченности цепи эти точки лежат в некотором $A \in L_\epsilon$ и, следовательно, $\rho(x, y) \geq \epsilon$.
    А значит и $A_\epsilon \in \Phi_\epsilon$ и цепь имеет мажоранту.
    Таким образом у нас выполнены условия леммы Цорна, поэтому в $\Phi_\epsilon$ имеется максимальный элемент $A_{\epsilon}^{\max}$.
    В силу максимальности у нас для всякого элемента вне $A^{\max}_{\epsilon}$ верно, что его расстояние до хотя бы какой-то точки из множества меньше $\epsilon$.
    Предположим, что у нас нет более чем счётного $\epsilon$-разреженного множества ни при каком $\epsilon$.
    Тогда максимальный элемент в $\Phi_\epsilon$ не более чем счётен при всяком $\epsilon$, и
    \[
        A = \bigcup_{n = 2}^{\infty} A^{\max}_{\frac{\epsilon_0}{n}}
    \]
    также не более чем счётно.
    Но, заметим, что $[A] = X$ -- а значит $X$ является сепарабельным -- противоречие.
    Поэтому при некотором $n$ множество $A^{\max}_{\frac{\epsilon_0}{n}}$ является более чем счётным -- что и требовалось найти.
\end{proof}
\begin{example}
    Рассмотрим $(l^\infty, \|\cdot\|_\infty)$.
    Пусть $\mathcal{B} = 2^{\N} \setminus \{\emptyset\}$.
    Рассмотрим
    \[
        S = \{\Ind_{A} \mid A \in \mathcal{B}\}.
    \]
    Это множество более чем счётно и вложено в $l^\infty$.
    Заметим, что оно является $1$-разреженным:
    \[
        \forall A, B \in \mathcal{B}\colon A \neq B \hookrightarrow \|\Ind_A - \Ind_B\|_{\infty} = \|\Ind_{A \triangle B}\|_{\infty} = 1,
    \]
    поскольку $A \triangle B \neq \emptyset$.

    Следовательно, по критерию несепарабельности, пространство $l^\infty$ не является сепарабельным.
\end{example}
\section{Принцип вложенных шаров.}
\begin{theorem}
    Пусть $(X, \rho)$ -- метрическое пространство.
    Следующие условия эквивалентны:
    \begin{enumerate}
        \item $X$ -- полно.
        \item Для всякой последовательности вложенных замкнутых шаров $\{B_{r_n}(x_n)\}_{n = 1}^{\infty}$ такой, что $r_n \rightarrow 0$, выполняется
        \[
            \bigcap_{n = 1}^{\infty} B_{r_n}(x_n) \neq \emptyset.
        \]
    \end{enumerate}
\end{theorem}
\begin{proof}
    Пусть $X$ -- полно и $\{B_{r_n}(x_n)\}$ -- последовательность вложенных шаров со стремящимися к 0 радиусами.
    Заметим, что
    \[
        \forall n \in \N \ \forall p \in \N \hookrightarrow \rho(x_{n + p}, x_n) \leq r_n.
    \]
    Поскольку $r_n \rightarrow 0$, то последовательность центров шаров -- фундаментальна и сходится к некоторой $x \in X$.
    В тоже время, для всякого $n \in \N$ и $m \geq n\colon$
    \[
        \rho(x_n, x_m) \leq r_n \rightarrow \rho(x_n, x) \leq r_n,
    \]
    а значит $x \in B_{r_n}(x_n)$.
    Поэтому $x \in \bigcap_{n = 1}^{\infty} B_{r_n}(x_n)$ и импликация показана.

    Покажем теперь обратную импликацию.
    Пусть $\{x_n\}$ -- фундаментальная последовательность.
    Мы можем построить подпоследовательность вида $\{x_{n_k}\}$ такую, что
    \[
        \rho(x_{n_k}, x_m) \leq 2^{-(k + 1)}, \forall m \geq n_{k}.
    \]
    Покажем, что в таком случае $B_{2^{-k}}(x_{n_k}) \supset B_{2^{-k - 1}}(x_{n_{k + 1}})$.

    Пусть $x \in B_{2^{-(k + 1)}}(x_{n_{k + 1}})$.
    Тогда $\rho(x, x_{n_k}) \leq \rho(x, x_{n_{k + 1}}) + \rho(x_{n_{k + 1}}, x_{n_k}) \leq 2^{-(k + 1)} + 2^{-(k + 1)} \leq 2^{-k}$, а значит $x \in B_{2^{-k}}(x_{n_k})$.
    Поскольку $2^{-k} \rightarrow 0$, то есть некоторая точка $x^* \in \bigcap_{k = 1}^{\infty} B_{2^{-k}}(x_{n_k})$.
    Она, очевидно, является пределом $x_{n_k}$ (так как $\rho(x, x_{n_k}) \leq 2^{-k}$).
    Но тогда, в силу фундаментальности, она является и пределом $x_n$, а значит $X$ -- полно.
\end{proof}
\begin{note}
    Заметим, что требование $r_n \rightarrow 0$ -- существенно.

    Пусть $(X, \rho) = (\N, \rho)$, где
    \[
        \rho(m, n) =
        \begin{cases}
            0, & m = n \\
            1 + \frac{1}{m + n}, & m \neq n.
        \end{cases}
    \]

    Это пространство, очевидно, полно, поскольку фундаментальны здесь только стационарные последовательности (если $m \neq n$, то $\rho(m, n) \geq 1$).

    С другой стороны, $B_{1 + \frac{1}{2m}}(m) = \{m, m + 1, \ldots\}$ и
    \[
        \bigcap_{m = 1}^{\infty} B_{1 + \frac{1}{2m}}(m) = \emptyset.
    \]
\end{note}
